\section{Introduction}
\label{sec:intro}

The proliferation of Internet of Things (IoT) edge devices and cyber-physical infrastructure has dramatically expanded the attack surface of contemporary local area networks~\cite{apruzzese2023role}. Machine Learning-based Network Intrusion Detection Systems (ML-NIDS) have demonstrated superior detection capabilities on benchmark telemetry datasets, frequently reporting accuracy exceeding $99\%$ in laboratory evaluations~\cite{ring2019survey}. Consequently, embedding neural intrusion detectors directly onto low-power microcontroller units (MCUs) at network edges—such as the Espressif ESP32—has emerged as a compelling paradigm for real-time, perimeter-adjacent traffic inspection with sub-millisecond threat mitigation~\cite{david2021tensorflow, banbury2021benchmarking}.

Despite these theoretical advantages, deep neural networks deployed for security inspection exhibit a critical vulnerability termed \emph{shortcut learning}~\cite{geirhos2020shortcut, arp2022dos}. Shortcut learning occurs when a model identifies non-causal, spurious correlations within its training distribution (e.g., protocol header artifacts, static port allocations, synthetic inter-packet arrival tempos, or client IP clustering) that predict malicious labels with high mathematical fidelity in-domain, but fail catastrophically when evaluated across operational network domains~\cite{sommer2010outside}. When evaluated zero-shot on an unobserved target network, these models suffer an acute performance drop known as \emph{cross-domain collapse}~\cite{sarhan2023standard}.

Simultaneously, edge deployment onto resource-constrained microcontrollers necessitates aggressive model compression~\cite{warden2019tinyml}. Post-Training Quantization (PTQ) from 32-bit floating-point (FP32) to 8-bit signed integer (INT8) arithmetic reduces memory footprints by $70\text{--}75\%$ and replaces power-intensive floating-point arithmetic with efficient integer multiply-accumulate (MAC) units~\cite{jacob2018quantization, gholami2022survey}. However, the intersection of aggressive low-precision quantization and shortcut learning remains fundamentally understudied in cybersecurity literature:
\begin{quote}
\emph{Does 8-bit integer post-training quantization worsen catastrophic cross-domain collapse by destroying subtle decision boundaries, or does the loss of precision act as an implicit regularizer that dampens fragile, non-causal shortcut features?}
\end{quote}

To answer this question, this paper presents \textbf{ZeroEdge-NIDS}, an empirical evaluation and software-simulated edge deployment platform. We train a multi-layer perceptron on the canonical CIC-IDS-2017 dataset~\cite{sharafaldin2018cicids} using a standardized 16-feature representation that drops known superficial shortcut artifacts (such as IP addresses, ports, sequence numbers, and TCP flags). We quantize the model to an integer-only INT8 execution graph and evaluate both the uncompressed FP32 baseline and the quantized INT8 variant zero-shot across all 257,673 flows of the UNSW-NB15 dataset~\cite{moustafa2015unsw} spanning nine distinct attack categories.

\subsection{Key Contributions}
The core contributions of this study are summarized as follows:
\begin{itemize}
    \item \textbf{Empirical Quantification of Cross-Domain Collapse:} We rigorously document the severe cross-domain generalization penalty of edge NIDS under zero-shot transfer, showing an accuracy collapse of $-61.92\%$ (dropping from $97.34\%$ in-domain to $35.42\%$ cross-domain).
    \item \textbf{Statistical Proof of Quantization Mitigation:} We demonstrate that INT8 post-training quantization provides a statistically significant mitigating effect on cross-domain collapse. The INT8 model achieves $36.51\%$ accuracy ($+1.09\%$ absolute boost over FP32, reducing collapse by $2.32\%$). A paired McNemar correctness test establishes this improvement as statistically significant ($\chi^2 = 282.89, p < 10^{-15}$), corroborated by a bootstrap $95\%$ confidence interval strictly above zero ($[+0.0138, +0.0152]$).
    \item \textbf{Shortcut Learning Diagnosis:} Through systematic permutation feature importance probes across both domains, we isolate the specific flow metrics (e.g., forward byte totals and flow duration) that operate as training shortcuts, confirming that precision truncation in INT8 dampens over-fitted gradient trajectories.
    \item \textbf{Transparent Dual-Tier Hardware Characterization:} We present a principled benchmark framework that strictly delineates between empirical host CPU execution timings (measured via \texttt{time.perf\_counter\_ns} with garbage collection disabled) and analytical ESP32 microcontroller projections ($5,504$ INT8 cycles, $4,160$ bytes SRAM, and $4.69\ \mu\text{J}$ per inference), accompanied by an exportable C header (\texttt{model\_data.h}).
    \item \textbf{Interactive Experiment Lab Platform:} We release an open-source evaluation suite featuring an interactive oscilloscope UI and an experiment engine that enforces dependency gating, preventing classification metrics from being reported on unlabeled telemetry.
\end{itemize}

The remainder of this paper is structured as follows: Section~\ref{sec:related} surveys related literature. Section~\ref{sec:architecture} formulates the mathematical model, feature mapping, and quantization arithmetic. Section~\ref{sec:implementation} details software implementation and hardware modeling. Section~\ref{sec:setup} outlines experimental design. Section~\ref{sec:results} presents quantitative findings. Section~\ref{sec:discussion} analyzes trade-offs and limitations. Section~\ref{sec:future} discusses future work, and Section~\ref{sec:conclusion} concludes the paper.
